\section{논의와 활용}
\label{sec:discussion}

\subsection{검증 절차는 학습 데이터 준비를 어떻게 바꿀 수 있는가?}

검증 절차는 기록 궤적이나 주석 내용을 교정하지 않고, 학습 전에 사용할 수 있는 증거의 범위를 유형화한다. 원천 CoC와 품질 관문, 도출 자차 움직임 응답, 반복 정책과 잠정 제한 시간, 연결 관계 및 장면 연속성 출처를 함께 보존하면 동일한 샘플을 단순 통과/탈락으로 압축하지 않고 후속 처리 근거와 함께 기록할 수 있다. 특히 \texttt{orphan\_response}는 수락된 트리거를 제거하면서 연결 응답을 원자적으로 제거하거나 재연결하지 않은 \emph{합성 연결 오류}이므로, 정상 또는 위험 궤적 레이블이 아니라 전처리 회귀 시험으로만 사용한다.

표~\ref{tab:training-policy}와 그림~\ref{fig:training-routing}은 이 구분을 학습 전 분류 정책으로 표현한다. 원천ㆍ품질ㆍ계약 근거가 함께 유지된 항목은 ``검증 조건 충족 항목'', 정책이나 탐지기 선택에 따라 판정이 달라지는 항목은 ``검토 후보'', 연결 관계 또는 장면 연속성 근거가 없는 항목은 ``출처 불충분''으로 둔다. 알려진 인위적 오류 사례는 검사기 회귀 시험에만 사용한다. 이 분류는 본 논문이 \emph{제안하는 데이터 준비 정책}이며, 실제 학습 성능이나 폐루프 행동을 개선하는지는 평가하지 않았다.

\begin{table}[t]
\centering
\bitablecaption{검증 증거에 따른 제안 학습 전 분류.}{Proposed pre-training routing according to verification evidence.}
\label{tab:training-policy}
\footnotesize
\begin{tabularx}{\columnwidth}{L{0.43\columnwidth}Y}
\toprule
Verification evidence & Proposed route \\
\midrule
Source, quality, and contract evidence preserved & Contract-satisfied item \\
Sensitive to policy or detector setting & Review candidate \\
Missing linkage or continuity evidence & Insufficient provenance \\
Known synthetic mutation & Checker-regression use only \\
\bottomrule
\end{tabularx}
\end{table}

\begin{figure}[t]
\centering
\begin{tikzpicture}[
  node distance=5mm,
  box/.style={draw, rounded corners=2pt, align=center, minimum height=8mm,
    text width=0.42\columnwidth, font=\scriptsize},
  decision/.style={draw, diamond, aspect=2.8, align=center, fill=neutral!8,
    inner sep=2pt, font=\footnotesize},
  keep/.style={box, fill=observed!10},
  review/.style={box, fill=assumed!10},
  provenance/.style={box, fill=derived!8},
  synthetic/.style={box, fill=neutral!8},
  flow/.style={-{Latex[length=2mm]}, thick}
]
\node[decision] (route) {Route by verification evidence};
\node[keep] (a) at (-0.235\columnwidth,-1.6) {Response observed under\\selected settings\\review retention after\\evidence check};
\node[review] (b) at (0.235\columnwidth,-1.6) {Sensitive to repeat policy,\\candidate deadline, or\\detector setting\\human review required};
\node[provenance] (c) at (-0.235\columnwidth,-3.0) {Insufficient CoC--response\\linkage or scene provenance\\do not concatenate\\into a long sequence};
\node[synthetic] (d) at (0.235\columnwidth,-3.0) {Known synthetic mutation\\checker-regression use only};
\draw[flow] (route.south) -- (a.north);
\draw[flow] (route.south) -- (b.north);
% Route the lower branches around the upper boxes so that their arrows do
% not pass through the text boxes for (a) and (b).
\draw[flow] (route.west) -- (-0.48\columnwidth,0) |- (c.west);
\draw[flow] (route.east) -- (0.48\columnwidth,0) |- (d.east);
\end{tikzpicture}

\bifigcaption{검증 증거의 제안 학습 전 분류. 각 분기는 물리적 안전/위험 레이블이 아니다.}{Evidence-based pre-training routing; no branch is a physical safety or risk label.}
\label{fig:training-routing}
\end{figure}

\subsection{왜 계약 통과와 검토 후보는 안전 레이블이 아닌가?}

계약 통과는 기록된 CoC 사건과 도출 자차 움직임 응답이 선택한 반복 의미론, 잠정 제한 시간 및 탐지기 아래에서 일관적이라는 뜻이다. 검토 후보는 같은 정책 아래 추가 확인이 필요하다는 뜻이다. 어느 판정도 상대 객체 거리ㆍ속도, 액추에이터 명령, 차량 동역학, 마찰 또는 폐루프 결과를 포함하지 않으므로 물리적 안전이나 위험을 직접 나타내지 않는다. 또한 $D=3$초는 법규나 차량 요구사항에서 도출된 임계값이 아니며, \texttt{scene-0028}과 \texttt{scene-0074}의 \preserve/\reset 판정 반전 및 27개 탐지기 설정의 민감도는 분류가 정책 의존적임을 보여준다. \preserve는 최초 사건 기준으로 지연 해석을 바꾸는 정책이지 유일한 정답이 아니다. 따라서 125+4건은 안전 데이터가 아니고, 다섯 건은 위험 데이터가 아니라 유형이 명시된 검토 후보이다. 별도 파서ㆍ탐지기의 13체인 플래그는 이 분류의 정량 근거에 포함하지 않는다.

이 경계는 학습 데이터 선택과 안전성 평가의 역할을 분리한다. 검증 절차는 시간ㆍ출처 계약이 명확한 항목을 보존하고 모호한 항목의 검토 사유를 남길 수 있지만, 학습용 정답의 사실성이나 제어 행동의 적절성을 새로 입증하지 않는다. 더 강한 판정에는 독립 궤적 검증, 객체ㆍ위험 연속성, 요구사항에서 도출한 기한, 폐루프 및 물리 증거가 추가되어야 한다.

\subsection{집계 일치 기록이 있어도 정형 관찰자는 왜 유용한가?}

\texttt{workshop-experiment-results.json} 집계 산출물은 별도 체인 구현의 비교 가능한 판정이 13/13 일치했다고 기록한다. 실행 가능한 기준선이나 항목별 비교 출력이 보존되지 않았으므로 이를 독립 구현 재현성으로 확대하지 않는다. 또한 공유 계약과 다른 파서ㆍ탐지기의 기본 동작 기록이므로 \uppaal의 수치적 우월성이나 중심 주장의 정량 근거도 아니다.

정형 관찰자의 남는 가치는 계산 자체보다 명세 구조와 추적 가능성에 있다. 5개 템플릿 \texttt{scene-0001} 모델, 2개 템플릿 분석 대상 집합 모델, \texttt{LoopSource}/\texttt{Monitor} 체인 모델은 서로 다른 목적을 가지며 하나의 통합 모델이 아니다. 각 모델군에서 반복 시계 정책ㆍ연결 관계ㆍ출처 플래그ㆍ장면 경계ㆍ교착ㆍ인위적 오류 사례 오라클을 명시하고 판정 이유를 사건, 정책 상태 및 반례 궤적으로 역추적할 수 있다. 다만 현재 입력은 기록된 결정적 궤적의 재생이므로 이 장점도 일반 환경 탐색, 제어 합성 또는 물리 안전성 증명으로 확대하지 않는다.

\subsection{연속 CoC 검증의 효용}

합성 변형에서 상태형 규칙만 모순으로 판정한 20건은 모두 현재 사건만으로는 판정할 수 없고 이전 사건의 의무 상태를 유지하도록 설계한 유형이었다. 이는 사건별 규칙이 지속 의무의 망각과 해제 후 잔존을 표현할 수 없다는 구조적 차이를 확인한다. 그러나 생성기가 상태형 판정을 사후조건으로 강제했으므로, 이 20건은 자연 자료에서 추가로 찾은 오류가 아니다. 제안 검사를 학습 전 품질 관문으로 실제 사용하려면, 자연 발생 모순을 포함한 독립 평가셋에서 정밀도와 재현율을 별도로 측정해야 한다.

UPPAAL의 역할은 전체 운전 행동을 생성하거나 안전성을 증명하는 것이 아니라, 선언한 계약 의미론에서 Python 판정을 상태 전이 수준으로 재현하는 것이다. 활성 정지 의무, 해제 증거와 후속 진행 사건이 어떤 순서로 충돌했는지를 실행 경로로 제시하므로 단순 오류 레이블보다 추적 가능한 근거를 제공한다. 다만 UPPAAL 코드는 같은 전이 규칙에서 생성되므로 0/7 일치는 독립 검증이 아니다. 여러 CoC 시나리오를 하나의 전역 운전 모델로 합치는 문제는 객체 연속성, 환경과 제어 모델이 추가로 필요하므로 본 논문의 범위에 포함하지 않는다.

선택한 자연 표본에서도 강한 효용 주장을 할 수 없다. 네 검토자의 텍스트 일관성 Fleiss $\kappa$가 0.027이었고 27개 장면 모두 적어도 한 항목에서 의견이 갈렸다. 합의 결과는 0/27 모순이었으므로 현재 자료는 자연 모순 검출 효과를 보여 주지 않는다. 이 결과는 계약 범주와 판정 지침 자체를 먼저 정교화하고 독립적인 자연 모순ㆍ비모순 사례를 확보해야 함을 시사한다.
